\documentclass[12pt,oneside]{article}
\usepackage[T1]{fontenc}
\usepackage{amsmath,amssymb,array,longtable}
\usepackage[english]{babel}
\usepackage[bookmarks=false,pdfauthor={Tomi Setsu},pdftitle={FunC v0.5.0 Design Proposal}]{hyperref}
\usepackage{fancyhdr}

\setlength{\headheight}{15.2pt}
\pagestyle{fancy}
\fancyhf{}
\fancyfoot[C]{\thepage}
\renewcommand{\headrulewidth}{0.5pt}
\fancyhead[C]{FunC v0.5.0 Design Proposal}

\title{FunC v0.5.0\\Ergonomic and Strongly Typed Design Proposal}
\author{Tomi Setsu}

\begin{document}
\maketitle

\begin{center}
\textbf{Document version:} v3 (Draft)\quad$\vert$\quad
\textbf{Supersedes:} v1 (initial proposal), v2 (10 Canonical Decisions locked)
\end{center}

\begin{abstract}
This document proposes a new source-level design target for FunC,
version \verb|0.5.0|. It starts from the currently implemented FunC
\verb|0.4.6| surface documented in \verb|doc/func.tex|, but reworks the
language around five principles:

\begin{enumerate}
\item strong domain types for TVM/TOS values such as typed cells,
\verb|Slice|, \verb|Builder|, \verb|Address|, and typed messages;
\item algebraic data types and exhaustive pattern matching;
\item explicit fallibility through \verb|Option| and \verb|Result|;
\item explicit mutability and explicit side effects;
\item module-, trait-, and generic-oriented reuse suitable for a standard
library and contract ecosystem.
\end{enumerate}

The language direction intentionally moves toward Rust-like ergonomics, but does
not attempt to become ``Rust on TVM''. In particular, it does not introduce
ownership, borrowing, lifetimes, macro metaprogramming, or a synchronous call
model that would conflict with TOS's asynchronous, actor-like execution.

This is a design proposal, not an implementation manifest. Its role is to
define a coherent, teachable next-generation FunC surface that remains faithful
to TVM/TOS semantics.
\end{abstract}

\section*{Introduction}
Current FunC is compact and close to TVM, but that same closeness makes it hard
to learn and easy to misuse. Several recurring pain points follow directly from
the \verb|0.4.6| language shape:

\begin{itemize}
\item domain values are not modeled as strongly as they should be;
\item boolean logic, parsing success, and protocol errors are too often
represented as unstructured integers;
\item mutation is hidden inside special syntax such as \verb|~method|;
\item exceptions and implicit runtime conventions are used where typed return
values would be easier to reason about;
\item code reuse relies too heavily on low-level built-ins and ad hoc include
patterns.
\end{itemize}

FunC \verb|0.5.0| addresses these issues by keeping the low-level power of TVM
available while moving ordinary authoring toward a safer and more legible
surface. The intended result is not ``higher level at any cost'', but a better
alignment between source code and the actual constraints of on-chain execution.

\clearpage
\tableofcontents

\clearpage
\section{Status and Positioning}
\subsection{Status}
This document is a \textbf{draft design proposal}. It defines the desired
surface language, type system, and standard-library direction for a future FunC
\verb|0.5.0| edition. It does not claim that the current compiler already
implements these features.

\subsection{Relationship to FunC 0.4.6}
FunC \verb|0.4.6| is the implementation baseline. FunC \verb|0.5.0| is the
language-design target. The design aims to preserve the following invariants:

\begin{itemize}
\item compilation still targets TVM/Fift-style low-level code;
\item TVM cells, slices, builders, continuations, and message actions remain
first-class concerns;
\item asynchronous cross-contract messaging remains the fundamental execution
model;
\item low-level escape hatches continue to exist for system code and stdlib
authors.
\end{itemize}

\subsection{What This Proposal Is Not}
FunC \verb|0.5.0| is \emph{not} intended to be:

\begin{itemize}
\item a clone of Solidity;
\item a full Rust compiler for TVM;
\item a language that hides gas, bounce, serialization, or asynchronous
delivery;
\item a language that requires ownership and lifetime reasoning for ordinary
contract code.
\end{itemize}

\section{First Principles}
\subsection{Execution Model Comes First}
The language must reflect the real execution model of TVM/TOS rather than a
familiar but misleading model imported from EVM or general-purpose systems
programming.

The relevant protocol facts are:

\begin{enumerate}
\item each contract owns isolated state;
\item cross-contract interaction happens through asynchronous messages;
\item a message may bounce, arrive later, or never arrive;
\item fees, reserve actions, and outgoing messages are explicit protocol
operations;
\item serialization into cells is fundamental rather than incidental.
\end{enumerate}

These facts impose direct source-language consequences:

\begin{enumerate}
\item contract-visible types must distinguish addresses, messages, coins, and
raw cells;
\item fallibility should be part of the type system instead of being encoded as
integer conventions;
\item mutation should be visible in the source rather than hidden in naming
conventions;
\item side effects should be visible in signatures rather than collapsed into a
single catch-all \verb|impure| marker;
\item low-level TVM intrinsics should mostly live in the standard library, not
in the parser.
\end{enumerate}

\subsection{Human Factors Matter}
FunC \verb|0.5.0| is optimized for the tasks that contract authors actually
perform:

\begin{itemize}
\item define message schemas;
\item parse inbound messages;
\item update persistent state;
\item send typed outbound messages;
\item model recoverable failures clearly;
\item reuse audited components without reading raw TVM stack code.
\end{itemize}

The language should therefore bias toward:

\begin{itemize}
\item readable declarations;
\item local reasoning;
\item exhaustiveness checks;
\item effect visibility;
\item error paths that are explicit and grep-able.
\end{itemize}

\subsection{Non-Goals}
To preserve approachability, FunC \verb|0.5.0| deliberately excludes several
features often associated with Rust:

\begin{itemize}
\item no ownership-and-lifetime borrow checker in the Rust sense;
\item no lifetimes;
\item no user-defined macros;
\item no trait objects as a required abstraction mechanism;
\item no \verb|async/await|;
\item no hidden heap allocation model.
\end{itemize}

At the same time, the language \emph{may} include a lightweight checker for
mutable-place overlap. Rejecting obviously dangerous expressions such as
``mutate the same slice twice in one expression'' is compatible with the goal of
remaining easy to learn; it does not require authors to reason about lifetimes
or ownership graphs.

\section{Headline Changes with Respect to 0.4.6}
\begin{longtable}{|p{0.28\textwidth}|p{0.31\textwidth}|p{0.31\textwidth}|}
\hline
\textbf{Concern} & \textbf{FunC 0.4.6} & \textbf{FunC 0.5.0 proposal} \\ \hline
Local binding &
\verb|var x = ...| or explicit type application syntax &
\verb|let x = ...|, \verb|let mut x = ...|, and ordinary typed bindings \\ \hline
Boolean values &
commonly encoded as \verb|int| &
dedicated \verb|bool| type with no implicit int-to-bool coercion \\ \hline
Parsing failures &
integer flags or exceptions &
\verb|Option<T>| and \verb|Result<T, E>| \\ \hline
Mutating methods &
\verb|~method| syntax &
ordinary method syntax with \verb|mut self| in the signature \\ \hline
Error handling &
\verb|throw_*| helpers and \verb|try/catch| &
typed \verb|Result| plus \verb|?| and \verb|match|; raw abort still exists for low-level code \\ \hline
Domain types &
\verb|int/cell/slice/builder/cont| plus conventions &
rich built-in domain types such as \verb|Cell<T>|, \verb|Slice|,
\verb|Builder|, \verb|Address|, \verb|Coins|, and
\verb|InternalContext<T>| \\ \hline
Control over impurity &
single \verb|impure| qualifier &
explicit effect sets such as \verb|effect(read, write, send)| \\ \hline
Message schemas &
manual op-code parsing and slice plumbing &
\verb|enum|, \verb|struct|, codec traits, and typed entrypoints \\ \hline
Serialization &
manual builder/slice plumbing for every protocol type &
\verb|Cell<T>|, \verb|to_cell|, \verb|from_cell|, \verb|load_any|, and
\verb|store_any| as ordinary typed APIs \\ \hline
Mutable alias safety &
mostly unchecked beyond surface syntax &
explicit mutable arguments plus a lightweight overlap checker for mutable places \\ \hline
Reuse &
include-heavy and built-in-heavy &
\verb|mod|, \verb|use|, generics, \verb|trait|, and \verb|impl| \\ \hline
Built-ins &
many parser-predefined names &
small core + stdlib \verb|extern "tvm"| declarations \\ \hline
\end{longtable}

\section{Design Goals}
\subsection{Primary Goals}
FunC \verb|0.5.0| has the following goals:

\begin{enumerate}
\item preserve TVM/TOS fidelity;
\item reduce incidental syntax complexity;
\item make domain concepts first-class in the type system;
\item make recoverable failure ordinary and typed;
\item make mutation and effects explicit;
\item make modules, traits, and generics the primary reuse mechanism;
\item keep low-level interop available for advanced code and compiler-adjacent
libraries.
\end{enumerate}

\subsection{Compatibility Goals}
The language should permit a staged migration from existing FunC contracts.
Legacy code should not need to be rewritten in one step. Instead, the compiler
should support an edition or compatibility mode in which old constructs are
accepted while new constructs are gradually adopted.

\section{Lexical Structure}
\subsection{Whitespace and Comments}
FunC \verb|0.5.0| adopts a more conventional comment system:

\begin{itemize}
\item \verb|//| introduces a line comment;
\item \verb|/* ... */| introduces a nested block comment.
\end{itemize}

For migration purposes, the legacy FunC comment forms \verb|;;| and
\verb|{- ... -}| may remain accepted in compatibility mode, but they are not
part of the preferred \verb|0.5.0| surface.

\subsection{Identifiers}
Ordinary identifiers use a conventional snake-case-oriented shape:

\begin{verbatim}
[A-Za-z_][A-Za-z0-9_]*
\end{verbatim}

Module-qualified names use \verb|::|:

\begin{verbatim}
std::slice::SliceReader
wallet::msg::Transfer
\end{verbatim}

Predicate names should use \verb|is_*| or \verb|has_*| rather than the old
\verb|name?| convention. Fallible operations should use \verb|try_*| or return
\verb|Result|.

\subsection{Literals}
The preferred literal forms are:

\begin{itemize}
\item integers: \verb|0|, \verb|17|, \verb|0xff|, \verb|42u32|,
\verb|-1i257|;
\item booleans: \verb|true|, \verb|false|;
\item strings: \verb|"hello"|;
\item byte strings: \verb|b"abc"|;
\item hexadecimal byte strings: \verb|hex"deadbeef"|;
\item addresses: \verb|addr"EQC..."| or \verb|addr"-1:abcd..."|.
\end{itemize}

The old string-suffix forms from FunC \verb|0.4.6| may remain available behind
compatibility gates, but new code should prefer explicit prefixes and typed
constructors.

\subsection{Keywords}
The core language reserves the following new keyword set:

\begin{verbatim}
mod use pub const type struct enum message trait impl fn extern
let mut static match if else while loop for in break continue
return where self Self as
effect pure unsafe
entry external internal bounce get
\end{verbatim}

Two of these (\verb|static|, \verb|where|) are reserved for future use: the
language grammar in Section~\ref{sec:grammar} does not currently consume them.
They are listed here so that user identifiers do not collide.

Compatibility mode may additionally reserve legacy FunC words such as
\verb|global|, \verb|asm|, \verb|impure|, \verb|inline|, and
\verb|method_id|.

\subsection{Attributes}
\label{sec:attributes}
FunC \verb|0.5.0| introduces an attribute syntax for edition selection and
module-level metadata. Inner attributes apply to the enclosing item (here, the
source file):

\begin{verbatim}
InnerAttribute ::= "#" "!" "[" AttrBody "]"
OuterAttribute ::= "#"     "[" AttrBody "]"
AttrBody       ::= identifier [ "=" (string-literal | integer-literal) ]
                 | identifier "(" [ AttrArg { "," AttrArg } ] ")"
AttrArg        ::= identifier [ "=" (string-literal | integer-literal) ]
\end{verbatim}

The only recognised inner attribute in the edition-\verb|0.5.0| surface is
\verb|edition|:

\begin{verbatim}
#![edition = "0.5.0"]
#![edition = "0.4-compat"]
\end{verbatim}

Unknown attributes are reserved and reject at parse time to leave room for
future extensions. Outer attributes (\verb|#[...]|) on individual items are
reserved but unused in this edition.

\section{Source Structure and Modules}
\subsection{Compilation Unit}
A source file is a sequence of items:

\begin{verbatim}
Item ::= ModuleDecl
       | UseDecl
       | ConstDecl
       | TypeAlias
       | StructDecl
       | EnumDecl
       | MessageDecl
       | TraitDecl
       | ImplDecl
       | FunctionDecl
\end{verbatim}

\subsection{Modules}
Modules organize contract code and the standard library:

\begin{verbatim}
mod wallet::state;
mod wallet::msg;
use std::result::Result;
use std::message::{InternalContext, SendOptions};
\end{verbatim}

The language keeps modules simple:

\begin{itemize}
\item there is no macro expansion model tied to modules;
\item visibility is controlled by \verb|pub|;
\item module-qualified names use \verb|::|;
\item \verb|use| imports names into local scope.
\end{itemize}

\subsection{Public API Surface}
Every item is private to its module unless marked \verb|pub|. This matters for
standard library design, contract libraries, and test-only helpers:

\begin{verbatim}
pub struct WalletState { ... }
pub enum WalletMsg { ... }
pub message Transfer(op = 0x0f8a7ea5) { ... }
pub fn recv_internal(...) -> Result<(), WalletError> effect(write, send) { ... }
\end{verbatim}

\section{Type System}
\label{sec:v050-types}
\subsection{Built-In Scalar Types}
The proposed scalar types are:

\begin{verbatim}
bool
i32  i64  i128  i256  i257
u8   u16  u32  u64  u128  u256
\end{verbatim}

\verb|i257| remains available because TVM integers are fundamentally 257-bit
signed values. For migration, \verb|int| may remain an alias of \verb|i257| in
compatibility mode, but new code should prefer explicit widths.

\subsection{Built-In Domain Types}
The following types become first-class built-ins:

\begin{verbatim}
Cell<T>
Slice
Builder
Cont
Address
StdAddress
Coins
Bytes
RawCell
RawSlice
Vec<T>
EitherRefOrInline<T>
\end{verbatim}

The role of these types is:

\begin{itemize}
\item \verb|Cell<T>| represents a typed cell reference carrying an expected
decoded payload type;
\item \verb|Slice| and \verb|Builder| model the TVM read/write cell domain;
\item \verb|Address| models a full on-chain address, not an arbitrary slice;
\item \verb|Coins| models currency amounts and documents intent better than raw
\verb|u128|;
\item \verb|Bytes| represents ordinary byte strings;
\item \verb|RawCell| and \verb|RawSlice| are low-level escape-hatch types for
unchecked code;
\item \verb|Vec<T>| is a sequence type with a fixed stack and wire
representation, usable at getter and ABI boundaries
(\S\ref{sec:vec-type});
\item \verb|EitherRefOrInline<T>| is a field-level codec that chooses inline
or ref storage deterministically (\S\ref{sec:either-ref-inline}).
\end{itemize}

\subsection{\texttt{Vec<T>}}
\label{sec:vec-type}
A \verb|Vec<T>| is a homogeneous sequence usable wherever the built-in
domain types are usable, including struct fields, function arguments,
function returns, and \verb|get fn| exports.

\paragraph{Stack representation.}
A \verb|Vec<T>| on the TVM stack is a single TVM \verb|Tuple| whose elements
are the sequence's items in order. TVM tuples natively hold up to 255
elements; a \verb|Vec<T>| longer than 255 is represented by a chain of tail
tuples (last element of the outer tuple is itself a \verb|Vec<T>| for the
remainder). The chain is deterministic: each segment holds exactly
\verb|min(255, remaining)| elements before continuing.

\paragraph{Wire representation.}
When serialised into a cell, a \verb|Vec<T>| writes:

\begin{enumerate}
\item \verb|length:uint16| --- the number of items, big-endian;
\item the items in order, each by \verb|T|'s codec;
\item if the encoded items overflow the current cell (bit budget or ref
budget), the encoder spills into a \emph{tail ref} whose referent is a
cell containing the remaining items, recursively.
\end{enumerate}

Overflow is determined by the same fit predicate used in
\S\ref{sec:cd-bodylayout}. The \verb|uint16| prefix caps the wire length at
65\,535 items; \verb|Vec| values longer than that cannot be serialised and
raise \verb|CellError::Overflow| at the codec boundary.

\paragraph{Getter returns.}
A \verb|get fn| returning \verb|Vec<T>| exports the stack-tuple form
directly to the caller. Indexers and wallets that consume the getter see a
TVM tuple and can iterate without invoking a cell decode.

\subsection{\texttt{EitherRefOrInline<T>}}
\label{sec:either-ref-inline}
This type represents a field whose TL-B form is \verb|field:(Either T ^T)|:
one bit of discriminant followed by either an inline encoding of \verb|T| or
a reference to a cell carrying \verb|T|.

\paragraph{Stack representation.}
An \verb|EitherRefOrInline<T>| on the TVM stack is a \verb|T| (the
inline/ref distinction is a codec-level concern, not a stack concern). The
compiler carries no runtime discriminant past the codec boundary.

\paragraph{Wire representation.}
The encoder writes:

\begin{enumerate}
\item one discriminant bit (\verb|0| = inline, \verb|1| = ref);
\item if inline: the encoded \verb|T| immediately after the bit;
\item if ref: a single cell reference whose referent is the encoded \verb|T|.
\end{enumerate}

The choice is deterministic and uses the same fit predicate as
\S\ref{sec:cd-bodylayout}, specialised to the current field cursor: inline
iff the encoded \verb|T| fits in the remaining bits and refs of the current
cell after consuming the discriminant bit. Two compliant compilers must
produce identical bytes.

\verb|EitherRefOrInline<T>| is distinct from \verb|Option<Cell<T>>|: the
former is mandatory (\verb|T| is always present), the latter is
(\verb|0| = absent, \verb|1| = present).

\subsection{Typed Cell References}
Typed cells are central to making TVM serialization tractable without hiding it:

\begin{verbatim}
pub struct WalletState {
    seqno: u32,
    owner: Address,
    plugins: Cell<PluginDict>?,
}
\end{verbatim}

\verb|Cell<T>| is assignable to \verb|RawCell| when raw interop is needed, but
ordinary contract code should prefer the typed form. This improves reviewability
for persistent state, lazy-loaded message refs, and nested TL-B objects.

\subsection{Compound Types}
Compound types are:

\begin{verbatim}
(T1, T2, ..., Tn)          // tuples
[T; N]                     // fixed-size arrays
T?                         // sugar for Option<T>
Option<T>
Result<T, E>
\end{verbatim}

Tuples remain useful because TVM naturally works with multiple stack results,
but user-defined \verb|struct| and \verb|enum| become the preferred way to
describe meaningful protocol data.

\subsection{Struct Types}
Structures are declared as:

\begin{verbatim}
pub struct WalletState {
    seqno: u32,
    public_key: u256,
    plugins: Dict<PluginKey, ()>,
}
\end{verbatim}

Struct fields are named, typed, and immutable by default. Mutation happens
through a mutable binding or a method with \verb|mut self|.

\subsection{Enum Types}
Enumerations are algebraic data types and are central to the proposal:

\begin{verbatim}
pub enum WalletMsg {
    Transfer {
        query_id: u64,
        amount: Coins,
        to: Address,
        response_to: Option<Address>,
    },
    InstallPlugin {
        plugin: Address,
    },
    RemovePlugin {
        plugin: Address,
    },
}
\end{verbatim}

Enums may also carry explicit discriminants when required by wire formats:

\begin{verbatim}
pub enum Op : u32 {
    Transfer = 0x0f8a7ea5,
    Burn     = 0x595f07bc,
}
\end{verbatim}

\subsection{Message Declarations}
Wire-level message bodies deserve a dedicated declaration form because they are
one of the dominant concepts in TVM/TOS contracts:

\begin{verbatim}
pub message Transfer(op = 0x0f8a7ea5) {
    query_id: u64,
    amount: Coins,
    to: Address,
    response_to: Address?,
}
\end{verbatim}

A \verb|message| declaration behaves like a named struct plus a derived message
codec. In particular:

\begin{itemize}
\item the opcode binding is part of the declaration rather than an external
convention;
\item the declaration derives the relevant cell/TL-B codec traits;
\item typed entrypoints may request \verb|InternalContext<Transfer>| directly;
\item higher-level enums may still group multiple message bodies when a contract
wants one sum type for dispatch.
\end{itemize}

\subsection{Option and Result}
Two standard prelude types are assumed. Their variant order is fixed by the
canonical tag assignment in Section~\ref{sec:cd-stack-enum}:

\begin{verbatim}
pub enum Option<T> {
    None,        // tag 0
    Some(T),     // tag 1
}

pub enum Result<T, E> {
    Ok(T),       // tag 0
    Err(E),      // tag 1
}
\end{verbatim}

Wire-format and stack-layout are fixed canonically:
\verb|Option<T>| is bit-compatible with TL-B \verb|Maybe X|;
\verb|Result<T, E>| uses a 1-bit tag. See
Sections~\ref{sec:cd-stack-enum} and \ref{sec:cd-wire-enum}.

These types are not library afterthoughts. They are core to ordinary contract
authoring:

\begin{itemize}
\item \verb|Option<T>| expresses absence without sentinels;
\item \verb|Result<T, E>| expresses recoverable failure without bare exceptions;
\item \verb|match| and \verb|?| make both types ergonomic in routine code.
\end{itemize}

For readability, \verb|T?| is accepted as surface sugar for
\verb|Option<T>|:

\begin{verbatim}
Address?
Cell<JettonWalletState>?
\end{verbatim}

Unlike languages that normalize everything around \verb|null|, the teaching path
still centers \verb|Some|/\verb|None|, \verb|match|, and \verb|if let|.

\subsection{Type Aliases}
Type aliases remain important for readability:

\begin{verbatim}
pub type QueryId = u64;
pub type PluginKey = (i32, u256);
pub type Ton = Coins;
\end{verbatim}

\subsection{Generics}
Generic functions, structs, enums, and traits use angle-bracket syntax:

\begin{verbatim}
pub struct Dict<K, V> { ... }

pub fn first<T>(items: (T, T)) -> T { ... }
\end{verbatim}

Unlike Rust, the design does not introduce lifetime parameters or ownership
bounds. Generic programming is deliberately restricted to type parameters and
trait constraints.

\section{Traits and Implementations}
\subsection{Trait Declarations}
Traits provide interface-oriented reuse:

\begin{verbatim}
pub trait TlbCodec {
    fn store(self, mut b: Builder) -> Result<Builder, CellError> pure;
    fn load(mut s: Slice) -> Result<(Self, Slice), ParseError> pure;
}
\end{verbatim}

Inside a \verb|trait| or \verb|impl| body, \verb|Self| refers to the concrete
type that implements (or is being defined by) the surrounding block. Outside
these blocks \verb|Self| is not in scope. A trait method whose signature does
not name \verb|self| is a \textbf{static associated function} and is called
with explicit type syntax, e.g.\ \verb|WalletMsg::load(cs)|. Methods with a
\verb|self| or \verb|mut self| receiver are dispatched through ordinary
method-call syntax \verb|x.method(...)|. The proposal does \emph{not} include
borrowed receivers; there are no \verb|\&self| or \verb|\&mut self| forms.

\subsection{Implementations}
Implementations attach methods to types:

\begin{verbatim}
impl WalletState {
    pub fn increment_seqno(mut self) -> Self pure {
        self.seqno += 1;
        self
    }
}

impl TlbCodec for WalletMsg {
    fn store(self, mut b: Builder) -> Result<Builder, CellError> pure { ... }
    fn load(mut s: Slice) -> Result<(Self, Slice), ParseError> pure { ... }
}
\end{verbatim}

\subsection{Why Traits Instead of Parser Magic}
One design goal of \verb|0.5.0| is to move functionality out of the parser and
into audited, inspectable library code. Traits support that directly:

\begin{itemize}
\item codecs can be trait-based;
\item message sending can depend on a trait bound rather than raw cells;
\item collection and builder helpers can live in the standard library;
\item test doubles and mock implementations become natural.
\end{itemize}

This does \emph{not} mean the compiler should be blind to protocol structure.
A small amount of protocol-aware lowering remains justified for things such as
entrypoint contexts, message opcodes, and automatic inline-vs-ref message-body
placement. The goal is to eliminate accidental parser magic, not to forbid every
useful built-in lowering.

\section{Bindings, Mutability, and Patterns}
\subsection{Immutable by Default}
All local bindings are immutable unless declared with \verb|mut|:

\begin{verbatim}
let seqno: u32 = state.seqno;
let mut cs: Slice = body;
\end{verbatim}

This replaces the old mixture of \verb|var|, explicit type application, and
hidden mutation through \verb|~method| calls.

\subsection{Assignments}
Only mutable places may be assigned to:

\begin{verbatim}
let mut state = load_state()?;
state.seqno = state.seqno + 1;
\end{verbatim}

\subsection{Pattern-Oriented Control Flow}
Patterns are not limited to \verb|match|. The language should also support
\verb|if let| and \verb|while let| because optional values and typed message
variants are ubiquitous in contract code:

\begin{verbatim}
if let Some(response_to) = msg.response_to {
    send(response_to, Ack { query_id: msg.query_id }, opts)?;
}

while let Some(item) = queue.pop_front() {
    process(item)?;
}
\end{verbatim}

\subsection{Destructuring}
Patterns work in \verb|let| bindings and in \verb|match|:

\begin{verbatim}
let (amount, response_to) = (msg.amount, msg.response_to);
let WalletMsg::Transfer { amount, to, .. } = msg;
\end{verbatim}

Pattern forms include:

\begin{itemize}
\item wildcard \verb|_|;
\item name bindings;
\item tuple patterns;
\item struct patterns;
\item enum variant patterns.
\end{itemize}

\subsection{Mutating Methods}
Methods may declare \verb|mut self|:

\begin{verbatim}
impl Slice {
    pub fn load_uint(mut self, bits: u16) -> Result<u256, ParseError> pure { ... }
}
\end{verbatim}

At the call site, this looks ordinary:

\begin{verbatim}
let mut cs = body;
let op = cs.load_uint(32)?;
\end{verbatim}

This replaces the special \verb|cs~load_uint(32)| syntax. The source clearly
states mutability in the binding and in the method signature, rather than
hiding it in punctuation.

\paragraph{Error-path state of mutable receivers.} When a \verb|mut self|
method returns \verb|Result<_, E>|, the canonical semantics are:

\begin{itemize}
\item on \verb|Ok|, the caller-side mutable place is rebound to the updated
receiver, and the \verb|Ok| payload is propagated;
\item on \verb|Err|, the caller-side mutable place is \emph{restored to its
pre-call state}. Any partial cursor advancement or field mutation inside
the callee is discarded.
\end{itemize}

This means \verb|cs.load_uint(32)?| followed by code that reuses \verb|cs| in
the \verb|Err| branch can rely on \verb|cs| still pointing at the original
position. The lowering in \verb|doc/func_v0.5.0-lowering.tex| \S4.3 must
be read with this rule: the \texttt{unconditional writeback} shown there is a
conceptual desugar step, and the emitted code must guard the writeback with a
tag check, or must use a shadow copy, so that the \verb|Err| case does not
commit intermediate state. Authors may explicitly commit partial mutations by
writing \verb|let new_cs = cs; new_cs.load_uint(32)?| in a non-receiver form.

\subsection{Mutable Arguments and Lightweight Safety Checks}
\label{sec:mut-args}
Methods with \verb|mut self| remain the idiomatic surface for builder/slice-like
APIs. Free functions --- and in particular helper functions that update a
contract's persistent state --- accept mutable places explicitly through a
\verb|mut| parameter qualifier:

\begin{verbatim}
pub fn load_var_uint(mut s: Slice, bits: u16) -> Result<u256, ParseError> pure { ... }

let mut cs = body;
let n = load_var_uint(mut cs, 5)?;
\end{verbatim}

\paragraph{Helper functions that mutate persistent state.}
This is the canonical pattern for splitting a contract's entry body into
smaller helpers that each mutate the protocol state. There is no \verb|&self|
or \verb|&mut self| borrow form in FunC \verb|0.5.0|; state passes through
\verb|mut| arguments by value-thread:

\begin{verbatim}
fn on_transfer(
    mut state: WalletState,
    sender: Address,
    req: Transfer,
) -> Result<(), WalletError> effect(read, send) { ... }

fn recv_internal(ctx: InternalContext<WalletMsg>) -> Result<(), WalletError> ... {
    let mut state = load_state()?;
    match ctx.body {
        WalletMsg::Transfer(req) => on_transfer(mut state, ctx.sender, req)?,
        ...
    }
    save_state(state)
}
\end{verbatim}

The lowering is identical to a method with \verb|mut self|: the helper takes
the receiver by value and returns it updated on \verb|Ok|, rolled back on
\verb|Err| (per the rule in \S\ref{sec:cd-stack-enum} and
\verb|doc/func_v0.5.0-lowering.tex| \S4.3).

This pattern resolves what earlier drafts listed as ``helper functions need
\verb|&mut self|''; no such syntax is required.

\paragraph{Overlap checks.}
This proposal adopts three lightweight rules:

\begin{itemize}
\item only mutable lvalues may be passed as \verb|mut| arguments;
\item the same mutable place may not be passed twice in one full expression;
\item parent and child places may not be mutated simultaneously in one full
expression.
\end{itemize}

These rules are intentionally narrow. They prevent obvious aliasing hazards
without introducing ownership, region inference, or lifetime syntax.

\section{Functions and Effect System}
\subsection{Function Syntax}
The general form is:

\begin{verbatim}
[pub] fn name(arg1: T1, mut arg2: T2, ...) -> R [effect(...)] Block
\end{verbatim}

Examples:

\begin{verbatim}
pub fn min(x: i257, y: i257) -> i257 pure { ... }
pub fn load_state() -> Result<WalletState, ParseError> effect(read) { ... }
pub fn load_uint(mut s: Slice, bits: u16) -> Result<u256, ParseError> pure { ... }
pub fn send_transfer(...) -> Result<(), SendError> effect(send) { ... }
\end{verbatim}

\subsection{Pure by Default}
A function with no effect annotation is \verb|pure|. It may compute, match, and
transform values, but it may not:

\begin{itemize}
\item read persistent state or chain context;
\item modify persistent state;
\item emit outgoing messages;
\item use raw unsafe TVM intrinsics.
\end{itemize}

\subsection{Effect Sets}
The primary effect vocabulary is:

\begin{verbatim}
effect(read)
effect(write)
effect(send)
\end{verbatim}

These effects mean:

\begin{itemize}
\item \verb|read|: may inspect persistent state, inbound context, config, time,
balance, or other chain environment values;
\item \verb|write|: may persist updated state or otherwise mutate contract-owned
storage;
\item \verb|send|: may emit action-phase effects such as outbound messages,
reserve operations, or code changes.
\end{itemize}

An effect set is cumulative: \verb|effect(write, send)| may both write state
and send messages.

\verb|unsafe| is \emph{not} a member of the effect set. It is a separate
top-level qualifier on a function that grants permission to use raw
\verb|extern "tvm"| intrinsics and other surface-unsafe operations. See the
grammar in Section~\ref{sec:grammar} for the precise distinction: a function
signature may carry exactly one \texttt{EffectQualifier}, chosen from
\verb|pure|, \verb|unsafe|, or \verb|effect(...)|.

Effect propagation is transitive and mandatory: if a function calls another
function whose effect set is $E$, the caller's effect set must be a superset of
$E$ (or the caller itself must be \verb|unsafe|). A \verb|pure| function has
$E = \emptyset$ and may only call other \verb|pure| functions. \verb|mut self|
and \verb|mut| arguments are \emph{not} effects; they express local state
threading and are compatible with \verb|pure|.

\subsection{Why Not Keep \texttt{impure}}
The old \verb|impure| marker is too coarse. It does not distinguish a getter
that reads block context from a function that mutates storage and sends three
messages. In \verb|0.5.0|, effect sets become part of ordinary interface review.

\subsection{Entrypoints}
Contract-facing entrypoints are explicit:

\begin{verbatim}
entry external fn recv_external(
    ctx: ExternalContext<SignedRequest>,
) -> Result<(), WalletError> effect(read, write, send) { ... }

entry internal fn recv_internal(
    ctx: InternalContext<WalletMsg>,
) -> Result<(), WalletError> effect(read, write, send) { ... }

entry bounce fn on_bounce(
    ctx: BounceContext<RawSlice>,
) -> Result<(), WalletError> effect(read, write) { ... }
\end{verbatim}

Getter functions use a dedicated modifier:

\begin{verbatim}
get fn seqno() -> Result<u32, ParseError> effect(read) { ... }
\end{verbatim}

The compiler maps these to the existing TVM/TOS conventions for internal,
external, bounced, and get-method execution.

The context types are not ordinary heap objects. They are compiler-recognized
views over real VM inputs. Source code should nevertheless treat them as normal,
typed fields:

\begin{verbatim}
ctx.sender
ctx.value
ctx.forward_fee
ctx.body
\end{verbatim}

This makes the execution model explicit while keeping raw TVM stack plumbing out
of routine contract code.

\subsection{Entry Dispatch Matrix}
\label{sec:entry-dispatch}
Unlike FunC \verb|0.4.6| which uses a single \verb|recv_internal| and expects
the contract to branch on the \verb|bounced| header bit, \verb|0.5.0| splits
internal-message delivery into two entry families:

\begin{itemize}
\item fresh internal messages are delivered to \verb|entry internal|;
\item bounced messages are delivered to \verb|entry bounce|;
\item external messages are delivered to \verb|entry external|.
\end{itemize}

If a contract defines \verb|entry internal| but no \verb|entry bounce|, an
inbound bounced message is silently dropped after consuming only the minimal
gas needed by the low-level wrapper --- no user-written handler is invoked.
Symmetric rules apply to undefined \verb|entry internal| (fresh internal
silently dropped) and undefined \verb|entry external| (external message
rejected at the wrapper). Contracts requiring the old single-entry pattern
must use a \verb|BounceContext<RawSlice>| form and manually re-dispatch.

\subsection{Entry Decoding Semantics}
\label{sec:entry-decode}
Typed body decoding happens in the generated wrapper, not in user code. The
canonical semantics for decode failure differ per entry kind:

\begin{itemize}
\item \verb|entry internal|: if the inbound body cannot be decoded as the
requested \verb|T|, the wrapper exits successfully without invoking the user
body. No state is written, no messages are sent. This preserves the TOS
convention that unrecognized ops do not cause bounce-refund gas drain.
\item \verb|entry external|: decoding is lazy. The wrapper exposes
\verb|ctx.body_slice: Slice| and a user-invoked
\verb|ctx.parse_body<T>() -> Result<T, ParseError>|.
The wrapper does \emph{not} decode eagerly,
because the external-message gas limit (typically 10\,000 gas before
\verb|accept_message|) may be exhausted by decoding a large body. Authors
must call \verb|accept_message| first, then parse.
\item \verb|entry bounce|: decoding uses the bounced-message body convention
(at most 256 bits after the \verb|0xFFFFFFFF| header). The default
\verb|BounceContext<RawSlice>| form accepts any bounce; typed variants must
opt in and must be robust to truncation.
\end{itemize}

\paragraph{Bounced body width rule.}
For any \verb|entry bounce fn ...(ctx: BounceContext<T>)| with \verb|T| not
equal to \verb|RawSlice|, the compiler must statically verify
\[
\verb|max_encoded_bits|(T) \;\le\; 256
\]
where \verb|max_encoded_bits| is the worst-case bit width of \verb|T|'s
derived codec, not counting refs. The \verb|0xFFFFFFFF| bounce header
prefix occupies 32 bits and is added by the wrapper, not counted against the
256-bit budget available to \verb|T|. Ref counts are also verified: \verb|T|
must encode as purely inline bits (no refs) because bounced bodies are
truncated, not referenced.

If \verb|max_encoded_bits(T) > 256| or \verb|T| requires any refs, the
compile error is \verb|BounceBodyTooWide|, and the author must either
narrow \verb|T|'s schema or fall back to \verb|BounceContext<RawSlice>| and
parse manually.

If a contract wants eager decoding in \verb|entry external|, it must write
\verb|entry external(eager) fn ...|. See
Section~\ref{sec:cd-external-eager} for the canonical rule.

\subsection{External Intrinsics}
Raw TVM primitives should be declared as \verb|extern "tvm"| in the standard
library rather than being magic parser names:

\begin{verbatim}
extern "tvm" fn send_raw_message(msg: RawCell, mode: u8) -> () effect(send, unsafe);
extern "tvm" fn accept_message() -> () effect(send, unsafe);
\end{verbatim}

This preserves power while making the surface language smaller and more
teachable.

\section{Expressions and Statements}
\subsection{Core Statements}
The language supports:

\begin{verbatim}
let
let mut
if / else
if let
match
while
while let
loop
for
break
continue
return
\end{verbatim}

\subsection{Match Expressions}
Pattern matching is exhaustive by default:

\begin{verbatim}
match msg {
    WalletMsg::Transfer { query_id, amount, to, response_to } => { ... }
    WalletMsg::InstallPlugin { plugin } => { ... }
    WalletMsg::RemovePlugin { plugin } => { ... }
}
\end{verbatim}

This is a major ergonomics improvement over manually parsing an op code into an
\verb|int| and then maintaining a parallel maze of \verb|if| branches.

\subsection{Flow-Sensitive Narrowing}
Pattern matching and optional-value tests should refine types in a
flow-sensitive way:

\begin{verbatim}
if let Some(reply_to) = msg.response_to {
    send(reply_to, Receipt { query_id: msg.query_id }, opts)?;
}
\end{verbatim}

Inside the \verb|if let| body, \verb|reply_to| is an \verb|Address| rather than
\verb|Address?|. Similar narrowing should happen for \verb|match| arms and other
obvious success paths over \verb|Option|, \verb|Result|, and message variants.

\subsection{The Question-Mark Operator}
If a function returns \verb|Result<T, E>|, the \verb|?| operator may be applied
to intermediate \verb|Result| values:

\begin{verbatim}
let mut cs = body;
let global_id = cs.load_int(32)?;
let seqno = cs.load_uint(32)?;
let msg = WalletMsg::load(cs)?;
\end{verbatim}

This drastically reduces the need for hand-written \verb|throw_unless|,
\verb|try/catch|, and integer-status plumbing in ordinary code.

\subsection{Loop Design}
The preferred loop forms are:

\begin{verbatim}
while cond { ... }
loop { ... }
for item in items { ... }
\end{verbatim}

Legacy \verb|repeat| and \verb|do ... until| may remain in compatibility mode,
but the core language should converge on a smaller, more conventional control
surface.

\paragraph{Iterator protocol for \texttt{for}.}
\verb|for item in items| desugars to a \verb|while let Some(item) = iter.next()|
loop where \verb|iter| is derived from \verb|items| by the \verb|Iter| trait:

\begin{verbatim}
pub trait Iter {
    type Item;
    fn next(mut self) -> Option<Self::Item> pure;
}

pub trait IntoIter {
    type Item;
    type IntoIterTy: Iter<Item = Self::Item>;
    fn into_iter(self) -> Self::IntoIterTy pure;
}
\end{verbatim}

The desugaring rule is:

\begin{verbatim}
for <pattern> in <expr> { <body> }

    ==>

let mut __it = <expr>.into_iter();
while let Some(<pattern>) = __it.next() { <body> }
\end{verbatim}

In edition \verb|0.5.0| the compiler must accept \verb|for| only over
iterables whose type implements \verb|IntoIter|. The standard library provides
initial implementations for ranges, fixed arrays, tuples of uniform element
type, and dictionary iteration; additional implementations may be added by
stdlib or user code. Trait resolution is static (monomorphized), matching the
lowering rule in \verb|doc/func_v0.5.0-lowering.tex| \S4.10.

\section{Messages, Serialization, and TVM-Specific APIs}
\subsection{Typed Message Contexts}
Messages should be represented as typed values rather than raw slices whenever
possible. The standard library is expected to define context types such as:

\begin{verbatim}
pub struct InternalContext<T> {
    sender: Address,
    value: Coins,
    forward_fee: Coins,
    body: T,
}

pub struct ExternalContext<T> {
    body: T,
    signature: Option<Bytes>,
}

pub struct BounceContext<T = RawSlice> {
    sender: Address,
    value: Coins,
    original_forward_fee: Coins,
    body: T,
}
\end{verbatim}

When an entrypoint requests \verb|InternalContext<T>| for a concrete message
type, the compiler or stdlib-generated wrapper is responsible for parsing the
inbound body into \verb|T|. Contracts may still opt into \verb|RawSlice| when
they need manual parsing.

\subsection{Typed Cell APIs}
Serialization should be both trait-driven and convenient in routine code.
Per Section~\ref{sec:cd-coherence}, the \verb|to_cell| convenience lives on
the \verb|CellEncode| trait itself as a default method, not as a blanket
inherent impl:

\begin{verbatim}
pub trait CellEncode {
    fn encode(self) -> Result<RawCell, CellError> pure;
    fn to_cell(self) -> Result<Cell<Self>, CellError> pure {
        self.encode().map(Cell::from_raw)
    }
}

impl<T: CellDecode> Cell<T> {
    pub fn load(self) -> Result<T, ParseError> pure;
    pub fn begin_parse(self) -> Slice pure;
}
\end{verbatim}

The \verb|impl<T: CellDecode> Cell<T> { ... }| form \emph{is} admissible
because \verb|Cell| is declared in the standard library (the trait's
originating crate). User code may not add its own blanket impls over foreign
types.

\subsection{Codec Traits}
Serialization and deserialization should be expressed through traits:

\begin{verbatim}
pub trait CellEncode {
    fn encode(self) -> Result<RawCell, CellError> pure;
}

pub trait CellDecode: Sized {
    fn decode(cell: RawCell) -> Result<Self, ParseError> pure;
}
\end{verbatim}

For wire protocols that rely on TL-B, a narrower \verb|TlbCodec| trait may
specialize the same idea.

The standard library should also offer high-frequency helpers modeled directly on
how contracts are written:

\begin{verbatim}
impl Slice {
    pub fn load_any<T: CellDecode>(mut self) -> Result<T, ParseError> pure;
}

impl Builder {
    pub fn store_any<T: CellEncode>(mut self, value: T) -> Result<Self, CellError> pure;
}
\end{verbatim}

\subsection{Slice and Builder APIs}
The low-level APIs remain essential, but their contracts become explicit:

\begin{verbatim}
impl Slice {
    pub fn load_uint(mut self, bits: u16) -> Result<u256, ParseError> pure;
    pub fn load_int(mut self, bits: u16) -> Result<i257, ParseError> pure;
    pub fn load_ref(mut self) -> Result<RawCell, ParseError> pure;
    pub fn load_addr(mut self) -> Result<Address, ParseError> pure;
    pub fn remaining_bits(self) -> u16 pure;
    pub fn remaining_refs(self) -> u8 pure;
}

impl Builder {
    pub fn store_uint(mut self, value: u256, bits: u16) -> Result<Self, CellError> pure;
    pub fn store_int(mut self, value: i257, bits: u16) -> Result<Self, CellError> pure;
    pub fn store_ref(mut self, cell: RawCell) -> Result<Self, CellError> pure;
    pub fn end_cell(self) -> Result<RawCell, CellError> pure;
}
\end{verbatim}

Note that parsing and building remain explicit operations. The language does
not try to make TVM cells disappear; it makes them typed and tractable.

\subsection{Sending Messages}
High-level message sending should take typed bodies:

\begin{verbatim}
pub enum BounceMode {
    NoBounce,
    Body256,
    Rich,
}

pub enum BodyLayout {
    Auto,
    Inline,
    ByRef,
}

pub struct SendOptions {
    value: Coins,
    bounce: BounceMode,
    body_layout: BodyLayout,
    mode: SendMode,
}

pub fn send<T: CellEncode>(
    dst: Address,
    body: T,
    opts: SendOptions,
) -> Result<(), SendError> effect(send);
\end{verbatim}

This API expresses a real protocol operation more directly than ad hoc raw-cell
construction plus \verb|SENDRAWMSG|.

When \verb|body_layout| is \verb|Auto|, the compiler \emph{must} choose between
inline and ref according to the canonical fit predicate specified in
Section~\ref{sec:cd-bodylayout}. The predicate is a deterministic function of
\verb|(env_bits, env_refs, body_bits, body_refs)| and is stable across
toolchains: two compliant compilers compiling the same source produce
identical wire bytes. Using ``may'' here instead of ``must'' would silently
break message hashing, fee estimation, and receiver-side parsing across
contracts built with different toolchains.

When \verb|body| is the unit value \verb|()|, the generated envelope encodes
``no body''\,: the TL-B \verb|Maybe| bit for \verb|body:(Maybe \ \^{}Cell)| is
set to \verb|0|. This is distinct from a zero-width inlined body or an empty
cell reference.

\section{Error Model}
\subsection{Recoverable and Unrecoverable Failure}
FunC \verb|0.5.0| distinguishes:

\begin{itemize}
\item recoverable failures, represented by \verb|Result<T, E>|;
\item unrecoverable VM aborts, represented by \verb|abort| or raw unsafe
intrinsics.
\end{itemize}

The default teaching path should strongly prefer typed recoverable failure.

\subsection{Domain Error Types}
Contracts are expected to define domain-specific error enums:

\begin{verbatim}
pub enum WalletError {
    InvalidSignature,
    WrongGlobalId,
    WrongSeqno,
    UnknownPlugin,
    Parse(ParseError),
    Send(SendError),
}
\end{verbatim}

\subsection{Bridging Legacy Exit Codes}
Compatibility layers may map legacy numeric exit codes to domain enums:

\begin{verbatim}
pub enum LegacyExitCode : u16 {
    InvalidSignature = 202,
    WrongSeqno = 33,
}
\end{verbatim}

This lets modern source code remain typed even when it interoperates with
legacy tooling or protocol conventions.

\section{Standard Library Direction}
\subsection{Small Core, Rich Stdlib}
The language core should stay small. Most functionality that is currently
special or parser-defined should instead move to the standard library:

\begin{itemize}
\item arithmetic helper functions;
\item slice and builder utilities;
\item dictionary abstractions;
\item message and context types;
\item codec traits and implementations;
\item contract-state helpers;
\item cryptographic intrinsics surfaced as \verb|extern "tvm"| functions.
\end{itemize}

Expected stdlib conveniences include:

\begin{itemize}
\item \verb|create_message| and \verb|send| wrappers over raw action lists;
\item typed state load/save helpers built on \verb|Cell<T>|;
\item \verb|load_any| / \verb|store_any| helpers for routine TL-B work;
\item message-body declarations with derived codecs.
\end{itemize}

\subsection{Traits as the Unit of Reuse}
Traits provide the right level of abstraction for:

\begin{itemize}
\item codecs;
\item formatting and debugging;
\item message sending;
\item dictionary key/value constraints;
\item reusable protocol behaviors.
\end{itemize}

This gives TOS a realistic path toward an audited contract standard library
without forcing all composition through includes and naming conventions.

\section{What Remains Explicit}
The proposal intentionally keeps several low-level concerns visible:

\begin{itemize}
\item serialization into cells;
\item bounce handling;
\item coins attached to outgoing messages;
\item effectful operations such as \verb|send| and reserve actions;
\item unsafe raw TVM interop where needed.
\end{itemize}

This is a language for TVM/TOS, not an attempt to hide the platform.

\section{Compatibility and Migration}
\subsection{Edition Model}
The recommended migration path is edition-based. A source file may opt into the
new syntax and semantics explicitly:

\begin{verbatim}
#![edition = "0.5.0"]
\end{verbatim}

Legacy files may continue to compile under a compatibility edition:

\begin{verbatim}
#![edition = "0.4-compat"]
\end{verbatim}

\subsection{Compatibility Surface}
The compatibility edition may continue to accept:

\begin{itemize}
\item legacy type names such as \verb|int| and \verb|cell|;
\item legacy comments;
\item legacy \verb|global|, \verb|const|, \verb|asm|, \verb|impure|,
\verb|inline|, \verb|method_id| forms;
\item legacy \verb|~method| mutation syntax;
\item legacy exception-style helpers.
\end{itemize}

\subsection{Illustrative Source Migration}
\begin{verbatim}
// 0.4.6 style
var cs = in_msg;
var signature = in_msg~load_bits(512);
throw_unless(33, msg_seqno == stored_seqno);

// 0.5.0 style
let mut cs = in_msg;
let signature = cs.load_bits(512)?;
if msg_seqno != stored_seqno {
    return Err(WalletError::WrongSeqno);
}
\end{verbatim}

\subsection{Migration Principle}
The important migration principle is semantic, not cosmetic:

\begin{itemize}
\item move from raw integers toward explicit types;
\item move from hidden mutation toward \verb|mut|;
\item move from integer-status parsing toward \verb|Result|;
\item move from raw op-code dispatch toward \verb|enum + match|;
\item move from parser magic toward stdlib declarations, while retaining a small
amount of protocol-aware lowering where it clearly improves source-level
clarity.
\end{itemize}

\section{Worked Comparison: One Wallet Contract in 0.4.6 and 0.5.0}
The following comparison shows the same piece of contract logic in both
language styles:

\begin{enumerate}
\item load wallet state;
\item parse an inbound transfer request;
\item validate the transfer sequence number;
\item send native coins to the destination;
\item increment and persist the sequence number.
\end{enumerate}

The examples are intentionally compact. They are not meant to be full
production wallets; they are meant to expose the source-language shape of the
same task.

\subsection{FunC 0.4.6 Style}
In current-style FunC, the code is dominated by raw slice parsing, implicit
mutation through \verb|~method| syntax, integer-coded failures, and manual
message construction:

\begin{verbatim}
;; FunC 0.4.6 style

(int, slice) load_data() inline {
  var ds = get_data().begin_parse();
  return (ds~load_uint(32), ds~load_msg_addr());
}

() save_data(int seqno, slice owner) impure inline {
  set_data(
    begin_cell()
      .store_uint(seqno, 32)
      .store_slice(owner)
      .end_cell()
  );
}

() recv_internal(slice in_msg_body) impure {
  var (stored_seqno, owner) = load_data();
  var cs = in_msg_body;

  int op = cs~load_uint(32);
  throw_unless(40, op == 0x0f8a7ea5);

  int msg_seqno = cs~load_uint(32);
  slice to = cs~load_msg_addr();
  int amount = cs~load_tomis();

  throw_unless(33, msg_seqno == stored_seqno);
  throw_unless(34, amount > 0);

  send_raw_message(
    begin_cell()
      .store_uint(0x18, 6)
      .store_slice(to)
      .store_tomis(amount)
      .store_uint(0, 1 + 4 + 4 + 64 + 32 + 1 + 1)
      .end_cell(),
    1
  );

  save_data(stored_seqno + 1, owner);
}
\end{verbatim}

\subsection{FunC 0.5.0 Style}
In the proposed \verb|0.5.0| surface, the same logic is expressed through typed
state, a typed message declaration, explicit mutability, explicit effects, and a
\verb|Result|-based error path:

\begin{verbatim}
#![edition = "0.5.0"]

pub struct WalletState {
    seqno: u32,
    owner: Address,
}

pub message Transfer(op = 0x0f8a7ea5) {
    seqno: u32,
    to: Address,
    amount: Coins,
}

pub enum WalletError {
    WrongSeqno,
    ZeroAmount,
    Parse(ParseError),
    Send(SendError),
}

pub fn load_state() -> Result<WalletState, ParseError> effect(read) { ... }
pub fn save_state(state: WalletState) -> Result<(), CellError> effect(write) { ... }

entry internal fn recv_internal(
    ctx: InternalContext<Transfer>,
) -> Result<(), WalletError> effect(read, write, send) {
    let mut state = load_state()?;
    let transfer = ctx.body;

    if transfer.seqno != state.seqno {
        return Err(WalletError::WrongSeqno);
    }
    if transfer.amount.is_zero() {
        return Err(WalletError::ZeroAmount);
    }

    send(
        transfer.to,
        (),
        SendOptions {
            value: transfer.amount,
            bounce: BounceMode::Body256,
            body_layout: BodyLayout::Auto,
            mode: SendMode::PayFeesSeparately,
        }
    )?;

    state.seqno += 1;
    save_state(state)?;
    Ok(())
}
\end{verbatim}

\subsection{What the Comparison Shows}
The comparison illustrates the intended source-level improvements:

\begin{itemize}
\item \textbf{Typed message declaration instead of raw op-code dispatch.}
In \verb|0.4.6| the contract manually loads \verb|op|, then loads each field
from a \verb|slice|. In \verb|0.5.0| the \verb|Transfer| message binds its
opcode in the declaration and the entrypoint receives a typed body directly.

\item \textbf{Explicit domain types.}
\verb|Address| and \verb|Coins| replace raw \verb|slice| and \verb|int|
conventions, making the meaning of each field visible at the type level.

\item \textbf{Explicit mutation.}
In \verb|0.4.6| mutation is hidden inside \verb|cs~load_uint(...)| and similar
forms. In \verb|0.5.0| mutation is carried by \verb|let mut cs| or
\verb|let mut state| and by explicit mutable parameters.

\item \textbf{Typed recoverable errors instead of numeric abort codes.}
\verb|throw_unless(33, ...)| and \verb|throw_unless(34, ...)| become
\verb|Err(WalletError::WrongSeqno)| and \verb|Err(WalletError::ZeroAmount)|.
This makes reviews and refactors much easier.

\item \textbf{High-level send API instead of manual action-cell assembly.}
The \verb|0.4.6| version has to build an outbound message cell explicitly.
The \verb|0.5.0| version exposes the protocol operation as \verb|send(...)|,
with value, bounce mode, body layout, and mode still explicit.

\item \textbf{Persistent state becomes a typed struct.}
The state layout is no longer ``a sequence of manually loaded fields'' from the
reader's point of view. The source speaks in terms of \verb|WalletState|.
\end{itemize}

\subsection{Why This Example Matters}
This example is deliberately small, but the same pattern scales to larger
contracts:

\begin{itemize}
\item Jetton and NFT messages become enums instead of raw op-code trees;
\item single message bodies can be declared with bound opcodes instead of
parallel ``struct + opcode constant'' conventions;
\item parser helpers return \verb|Result| instead of integer flags;
\item bounce handlers match typed variants instead of re-parsing raw slices;
\item audited stdlib modules provide \verb|send|, codecs, and dict abstractions
instead of every contract open-coding them.
\end{itemize}

\section{Core Grammar Summary}
\label{sec:grammar}
The grammar below is intentionally schematic. It describes the shape of the
language, not every precedence rule or parsing ambiguity.

\begin{verbatim}
SourceFile ::= { InnerAttribute } { Item }

Item ::= { OuterAttribute }
         ( ModuleDecl
         | UseDecl
         | ConstDecl
         | TypeAlias
         | StructDecl
         | EnumDecl
         | MessageDecl
         | TraitDecl
         | ImplDecl
         | FunctionDecl )

ModuleDecl ::= "mod" Path ";"
UseDecl    ::= "use" Path [ "as" identifier ] ";"
ConstDecl  ::= [ "pub" ] "const" identifier ":" Type "=" Expr ";"
TypeAlias  ::= [ "pub" ] "type" identifier [ GenericParams ] "=" Type ";"

StructDecl ::= [ "pub" ] "struct" identifier [ GenericParams ] "{" Fields "}"
EnumDecl   ::= [ "pub" ] "enum" identifier [ GenericParams ] [ ":" Type ] "{"
               Variants "}"
MessageDecl ::= [ "pub" ] "message" identifier [ "(" "op" "=" Expr ")" ]
                "{" Fields "}"
TraitDecl  ::= [ "pub" ] "trait" identifier [ GenericParams ] "{"
               TraitItems "}"
ImplDecl   ::= "impl" [ GenericParams ] Type [ "for" Type ] "{"
               ImplItems "}"

FunctionDecl ::= [ "pub" ] [ EntryQualifier ] "fn" identifier
                 [ GenericParams ] "(" Params ")" "->" Type
                 [ EffectQualifier ] BlockOrSemicolon

EntryQualifier ::= "entry" ( "external" | "internal" | "bounce" )
                 | "get"

EffectQualifier ::= "pure"
                  | "unsafe"
                  | "effect" "(" EffectList ")"

Stmt ::= LetStmt
       | ExprStmt
       | ReturnStmt
       | IfStmt
       | IfLetStmt
       | MatchStmt
       | WhileStmt
       | WhileLetStmt
       | LoopStmt
       | ForStmt
       | BreakStmt
       | ContinueStmt
       | Block

Type ::= SimpleType
       | Type "?"

Arg ::= Expr
      | "mut" Expr
\end{verbatim}

\section{Illustrative Contract Fragment}
The following fragment demonstrates the overall style target:

\begin{verbatim}
#![edition = "0.5.0"]

use std::result::Result;
use std::message::{InternalContext, SendOptions, SendMode, send};

pub struct WalletState {
    seqno: u32,
    owner: Address,
}

pub enum WalletError {
    WrongSeqno,
    Parse(ParseError),
    Send(SendError),
}

pub message Transfer(op = 0x0f8a7ea5) {
    query_id: u64,
    amount: Coins,
    to: Address,
    response_to: Address?,
}

pub fn load_state() -> Result<WalletState, ParseError> effect(read) { ... }
pub fn save_state(state: WalletState) -> Result<(), CellError> effect(write) { ... }

entry internal fn recv_internal(
    ctx: InternalContext<Transfer>,
) -> Result<(), WalletError> effect(read, write, send) {
    let mut state = load_state()?;
    let msg = ctx.body;

    if let Some(response_to) = msg.response_to {
        send(
            response_to,
            Receipt { query_id: msg.query_id },
            SendOptions {
                value: 0.into(),
                bounce: BounceMode::NoBounce,
                body_layout: BodyLayout::Auto,
                mode: SendMode::PayFeesSeparately,
            },
        )?;
    }

    send(
        msg.to,
        msg,
        SendOptions {
            value: msg.amount,
            bounce: BounceMode::Body256,
            body_layout: BodyLayout::Auto,
            mode: SendMode::PayFeesSeparately,
        },
    )?;

    state.seqno = state.seqno + 1;
    save_state(state)?;
    Ok(())
}
\end{verbatim}

\section{Conclusion}
FunC \verb|0.5.0| should become easier to read not by hiding TVM/TOS, but by
modeling it honestly:

\begin{itemize}
\item typed messages instead of raw slices;
\item enums and exhaustive \verb|match| instead of integer dispatch;
\item message declarations with bound opcodes instead of ad hoc constants;
\item typed cell references and auto codec helpers instead of repetitive
builder/slice plumbing;
\item \verb|Result| and \verb|Option| instead of status flags and surprise
exceptions;
\item \verb|let mut| and \verb|mut self| instead of punctuation-based mutation;
\item \verb|if let| and flow-sensitive narrowing for common optional/message
control flow;
\item explicit effects instead of a single opaque impurity bit;
\item trait- and module-based reuse instead of parser magic and ad hoc
includes.
\end{itemize}

That combination preserves the strengths of FunC --- closeness to TVM, direct
control over serialization and messaging, and efficient compilation --- while
making the source language much easier to teach, review, and scale across a
real contract ecosystem.

\section{Canonical Decisions}
\label{sec:canonical-decisions}
The items below were open questions in an earlier draft. Each has now been
fixed from first principles: minimum runtime cost, determinism across
toolchains, fidelity to TVM/TOS, and alignment with existing TON conventions.
Compilers claiming specification conformance must obey these rules exactly.

\subsection{Canonical Decision 1: Stack Representation of Sum Types}
\label{sec:cd-stack-enum}
An \verb|enum|, \verb|Option<T>|, or \verb|Result<T, E>| value occupies
exactly two TVM stack values:

\begin{verbatim}
(tag : int,  payload : tuple)
\end{verbatim}

\begin{itemize}
\item \verb|tag| is an ordinary TVM 257-bit signed integer. No bit packing is
performed on the stack; that is a wire-format concern only.
\item \verb|payload| is an ordinary TVM tuple. Variant fields appear in
declaration order. A zero-field variant uses the empty tuple \verb|[]|.
A one-field variant uses the singleton tuple \verb|[v]|. Multi-field
variants use tuples of the appropriate arity.
\end{itemize}

Tag values:

\begin{itemize}
\item \verb|Option<T>|: \verb|None = 0|, \verb|Some = 1|.
\item \verb|Result<T, E>|: \verb|Ok = 0|, \verb|Err = 1|.
\item User-defined enums without an explicit discriminant type: tags are
assigned \verb|0, 1, 2, ...| in declaration order.
\item User-defined enums with an explicit discriminant type
(\verb|enum E : uN|) or explicit per-variant values: the declared numeric
values are the stack tags.
\end{itemize}

Optimisers may locally unbox the payload when the unboxing is not observable
at function, trait, or ABI boundaries, but the canonical model above is what
every compiler must be able to produce on demand (for example at
\verb|store_any| boundaries).

\subsection{Canonical Decision 2: Wire Format of Sum Types}
\label{sec:cd-wire-enum}
The codec for a sum type $S$ with $N$ variants writes:

\begin{enumerate}
\item a tag prefix of width
\[
w_S \;=\; \max\!\left(1,\; \lceil \log_2 N \rceil\right)
\]
bits, big-endian, unsigned, unless the enum declares an explicit discriminant
type \verb|: uN| in which case the tag is exactly \verb|uN| wide;
\item the variant fields, in declaration order, each by its own codec.
\end{enumerate}

For \verb|Option<T>| the tag is 1 bit (\verb|None = 0|, \verb|Some = 1|),
yielding a wire format bit-compatible with TL-B \verb|Maybe X|.
For \verb|Result<T, E>| the tag is 1 bit (\verb|Ok = 0|, \verb|Err = 1|).

A \verb|message M(op = X)| declaration is a special case: the wire prefix is
the 32-bit opcode \verb|X|, not a derived enum tag. A sum of \verb|message|
variants (for dispatch) uses the 32-bit opcodes of each constituent message
as its discriminant; the enum does not insert an additional tag.

Adding, removing, or reordering variants of an enum \emph{breaks the wire
format}. Implementations must reject silent breakage at compile time by
computing a schema hash over the enum definition and including it in the
generated ABI manifest.

\subsection{Canonical Decision 3: Layout of \texttt{Bytes}}
\label{sec:cd-bytes}
\verb|Bytes| is a stdlib struct with a fixed ABI:

\begin{itemize}
\item \textbf{Stack:} a single TVM \verb|Slice| that is guaranteed
byte-aligned (\verb|remaining_bits() % 8 == 0|) when entering any function
that accepts \verb|Bytes|.
\item \textbf{Wire:} the ``snake'' cell-chain convention already standard in
TON: up to $\lfloor (1023 - \mathit{header\,bits}) / 8 \rfloor$ bytes inline
in the current cell, overflow spilled into a single tail reference whose
referent is itself a snake cell. A length prefix is \emph{not} part of the
snake convention; the reader stops at end-of-slice.
\item For protocols that need a length-prefixed byte string, stdlib provides
\verb|LenPrefixedBytes| with an explicit \verb|u16| length.
\end{itemize}

\verb|Bytes| is deliberately not a \verb|Cell<Bytes>| because the most common
pattern (message bodies, metadata strings) inlines the prefix into the
current cell rather than forcing a ref.

\subsection{Canonical Decision 4: \texttt{BodyLayout::Auto} Fit Predicate}
\label{sec:cd-bodylayout}
Let \verb|env| be the current envelope builder after the compiler has stored
the common message header (source, destination, value, flags, created-lt,
created-at, init, and the outer \verb|Either| tag for the body). Let
\verb|body| be the encoded body after materialising it into a temporary
builder. Define:

\begin{verbatim}
env_bits   = env.bits_used()
env_refs   = env.refs_used()
body_bits  = body.bits_used()
body_refs  = body.refs_used()
\end{verbatim}

\verb|BodyLayout::Auto| emits the \emph{inline} form
(\verb|body:(Either X \^{}X) = Left X|) iff both of:

\begin{verbatim}
env_bits + 1 + body_bits <= 1023
env_refs +     body_refs <=    4
\end{verbatim}

Otherwise it emits the \emph{by-ref} form (\verb|Right \^{}X|). The constant
\verb|+1| accounts for the outer \verb|Either| tag bit. All compilers must
compute exactly this predicate.

When \verb|body_bits| is statically known (constant-size codec), the compiler
may fold the branch at compile time. When it is not, the compiler emits a
runtime branch whose test matches the predicate above bit-for-bit.

\verb|BodyLayout::Inline| and \verb|BodyLayout::ByRef| force the respective
choice; a body that does not fit inline under \verb|Inline| is a runtime
error mapped to \verb|CellError::Overflow|.

\subsection{Canonical Decision 5: Typed \texttt{Err} to Exit Code}
\label{sec:cd-err-exit}
When the user body of an \verb|entry| function returns \verb|Err(e)|:

\begin{itemize}
\item persistent state is \textbf{not} committed (no \verb|set_data|);
\item queued action-phase effects are \textbf{not} committed (the action list
built during the user body is discarded);
\item the contract terminates via \verb|THROW n| where \verb|n| is the
\emph{exit code} derived from \verb|e|.
\end{itemize}

The exit code derivation is:

\begin{enumerate}
\item If the error type is an \verb|enum|, and each variant carries an
explicit unsigned discriminant within \verb|[1, 65535]| (directly via
\verb|enum E : uN| and per-variant \verb|= N|), the discriminant is the
exit code.
\item Otherwise the compiler assigns
\verb|exit_code(variant_i) = 1000 + i| where \verb|i| is the zero-based
declaration index. The range \verb|[1000, 65535]| is the user range.
\item The range \verb|[0, 999]| is reserved: \verb|0| is success, \verb|1|
through \verb|999| are reserved for TVM, TOS, and compiler-internal
use. Explicit user discriminants in this range are a compile-time error.
\end{enumerate}

Variants that carry additional payload data (e.g.\ \verb|Parse(ParseError)|)
use only their own variant discriminant as the exit code; the nested payload
is discarded at the exit boundary. Contracts that need richer error
telemetry must emit it explicitly through \verb|send| or getters before
returning \verb|Err|.

On \verb|Ok|, the wrapper commits state (if the user called \verb|save_state|)
and action-phase effects normally, then returns \verb|0|.

\subsection{Canonical Decision 6: Getter / \texttt{method\_id} Selector Rule}
\label{sec:cd-method-id}
For every \verb|get fn name(...)|:

\begin{itemize}
\item the default exported method id is
\verb|crc16(name) \| 0x10000| (identical to FunC \verb|0.4.6|, preserving
interop with every existing TON/TOS indexer and wallet);
\item an explicit override uses the same syntax family as
\verb|0.4.6|: \verb|get fn name(...) method_id(N)| for a literal integer, or
\verb|method_id("string")| to select via the crc16 of an alternative name;
\item collisions in the final exported method-id set are a compile-time
error; the compiler reports every colliding pair;
\item two sibling compilations that produce different method-id assignments
for the same source are a violation of this rule.
\end{itemize}

The grammar for \verb|FunctionDecl| is extended accordingly:

\begin{verbatim}
FunctionDecl ::= [ "pub" ] [ EntryQualifier ] "fn" identifier
                 [ GenericParams ] "(" Params ")" "->" Type
                 [ EffectQualifier ]
                 [ "method_id" ( "(" IntLiteral ")"
                               | "(" StringLiteral ")" ) ]
                 BlockOrSemicolon
\end{verbatim}

\verb|method_id| is admissible only on \verb|get fn| functions in edition
\verb|0.5.0|; using it on a non-getter is a compile-time error.

\subsection{Canonical Decision 7: Expression Grammar and Precedence}
\label{sec:cd-expr}
The expression grammar below supersedes the earlier schematic summary. It is
parsed top-down with precedence climbing.

\begin{verbatim}
Expr       ::= Assign
Assign     ::= Or [ AssignOp Assign ]           // right-assoc
AssignOp   ::= "=" | "+=" | "-=" | "*=" | "/=" | "%="
             | "&=" | "|=" | "^=" | "<<=" | ">>="

Or         ::= And { "||" And }                  // left-assoc, short-circuit
And        ::= Compare { "&&" Compare }          // left-assoc, short-circuit
Compare    ::= BitOr [ CmpOp BitOr ]             // non-associative
CmpOp      ::= "==" | "!=" | "<" | ">" | "<=" | ">=" | "<=>"

BitOr      ::= BitXor { "|" BitXor }             // left-assoc
BitXor     ::= BitAnd { "^" BitAnd }             // left-assoc
BitAnd     ::= Shift  { "&" Shift }              // left-assoc
Shift      ::= Add    { ( "<<" | ">>" ) Add }    // left-assoc
Add        ::= Mul    { ( "+" | "-" ) Mul }      // left-assoc
Mul        ::= Unary  { ( "*" | "/" | "%" ) Unary }

Unary      ::= { UnaryOp } Cast
UnaryOp    ::= "-" | "!" | "~"
Cast       ::= Postfix [ "as" Type ]

Postfix    ::= Primary { PostfixOp }
PostfixOp  ::= "?"                               // Result/Option propagation
             | "." Ident [ "(" Args ")" ]        // field or method call
             | "(" Args ")"                      // function call
             | "[" Expr "]"                      // index
             | "::" Ident                        // associated fn/const path

Primary    ::= Literal
             | "(" ")" | "(" Expr { "," Expr } ")"
             | "[" "]" | "[" Expr { "," Expr } "]"
             | StructLit
             | EnumLit
             | Path
             | "if" Expr Block [ "else" Block ]
             | "match" Expr "{" { MatchArm } "}"
             | Block

Args       ::= [ Arg { "," Arg } ]
Arg        ::= Expr | "mut" Expr
\end{verbatim}

Precedence summary (highest to lowest):

\begin{center}
\begin{tabular}{|c|l|l|}
\hline
\textbf{Level} & \textbf{Forms} & \textbf{Associativity} \\ \hline
14 & postfix: \verb|?|, \verb|.|, \verb|(..)|, \verb|[..]|, \verb|::| & left \\ \hline
13 & prefix: \verb|-|, \verb|!|, \verb|~|; \verb|as| & right \\ \hline
12 & \verb|*|  \verb|/|  \verb|%| & left \\ \hline
11 & \verb|+|  \verb|-| & left \\ \hline
10 & \verb|<<|  \verb|>>| & left \\ \hline
 9 & \verb|&| & left \\ \hline
 8 & \verb|^| & left \\ \hline
 7 & \verb+|+ & left \\ \hline
 6 & \verb|==| \verb|!=| \verb|<| \verb|>| \verb|<=| \verb|>=| \verb|<=>| & non-associative \\ \hline
 5 & \verb|&&| & left (short-circuit) \\ \hline
 4 & \verb+||+ & left (short-circuit) \\ \hline
 2 & \verb|=| and compound assignments & right \\ \hline
\end{tabular}
\end{center}

Notes:

\begin{itemize}
\item Comparisons are non-chainable: \verb|a < b < c| is a syntax error.
\item \verb|?| binds tighter than method-call dot: \verb|x?.foo()| means
\verb|(x?).foo()|. Writing \verb|x.foo()?| propagates after calling.
\item \verb|as| takes a type, not an expression, on the right.
\item Bitwise \verb|&|/\verb+|+/\verb|^| coexist with logical
\verb|&&|/\verb+||+ at separate precedences. Logical operators require
\verb|bool| operands and short-circuit; bitwise operators work on
integer types only.
\end{itemize}

\subsection{Canonical Decision 8: Trait Coherence}
\label{sec:cd-coherence}
FunC \verb|0.5.0| adopts an orphan-style coherence rule:

\begin{itemize}
\item \verb|impl Trait for T| is admissible iff \verb|Trait| \emph{or}
\verb|T| is declared in the current compilation unit (the file or the
set of files forming one \verb|mod|).
\item A blanket \verb|impl<X: Bound> Trait for X| is admissible iff
\verb|Trait| is declared in the current compilation unit.
\item For any concrete pair \verb|(Trait, T)| there may be at most one
\verb|impl| in the full dependency graph; a second impl is a
compile-time error. This includes indirect duplication through blanket
impls.
\item The earlier \verb|impl<T: CellEncode> T { fn to_cell(...) }| shorthand
in \S\ref{sec:v050-types} is not admissible. \verb|to_cell| is instead a
method on the \verb|CellEncode| trait itself, either abstract or as a
default method using the trait's \verb|encode|:

\begin{verbatim}
pub trait CellEncode {
    fn encode(self) -> Result<RawCell, CellError> pure;
    fn to_cell(self) -> Result<Cell<Self>, CellError> pure {
        self.encode().map(Cell::from_raw)
    }
}
\end{verbatim}
\end{itemize}

Coherence is checked after monomorphisation: an instantiated concrete pair
must not acquire two impls via different routes.

\subsection{Canonical Decision 9: Eager Decoding on \texttt{entry external}}
\label{sec:cd-external-eager}
By default \verb|entry external fn foo(ctx: ExternalContext)| receives an
untyped \verb|ctx|:

\begin{verbatim}
pub struct ExternalContext {
    body_slice: Slice,
    signature: Option<Bytes>,
}

impl ExternalContext {
    pub fn parse_body<T: CellDecode>(mut self)
        -> Result<T, ParseError> pure { ... }
}
\end{verbatim}

Authors call \verb|parse_body| after \verb|accept_message|, so decoding gas
is paid from contract balance rather than from the pre-accept 10\,000-gas
external-message quota.

Authors who accept the gas risk may opt into eager decoding:

\begin{verbatim}
entry external(eager) fn recv_external(
    ctx: ExternalContext<SignedRequest>,
) -> Result<(), WalletError> effect(read, write, send) { ... }
\end{verbatim}

The grammar for \verb|EntryQualifier| is extended:

\begin{verbatim}
EntryQualifier ::= "entry" "external" [ "(" "eager" ")" ]
                 | "entry" "internal"
                 | "entry" "bounce"
                 | "get"
\end{verbatim}

For \verb|entry internal| and \verb|entry bounce| the default is already
eager decoding; no opt-in is required.

\subsection{Canonical Decision 10: Compatibility Edition Surface}
\label{sec:cd-compat}
An edition is chosen per source file by an inner attribute; it is not
per-item.

\begin{itemize}
\item \verb|#![edition = "0.5.0"]| selects the strict \verb|0.5.0| surface.
This is the recommended edition for new code.
\item \verb|#![edition = "0.4-compat"]| selects the superset edition that
accepts all \verb|0.4.6| constructs plus all \verb|0.5.0| constructs.
The declared purpose is migration, not long-term authoring.
\item Absence of an edition attribute defaults to \verb|0.4-compat| for
this release, so existing contracts continue to compile without
modification.
\end{itemize}

Under \verb|0.4-compat| the following legacy constructs are accepted in
addition to the \verb|0.5.0| surface:

\begin{itemize}
\item lexical: \verb|;;| line comments; \verb|{- ... -}| nested block
comments; legacy string suffixes \verb|s|, \verb|a|, \verb|u|, \verb|h|,
\verb|H|, \verb|c|; \verb|?|-suffixed predicate identifiers;
\verb|~|-prefixed modifying methods;
\item top-level: \verb|#pragma ...|, \verb|#include ...|, \verb|global|
declarations, \verb|const| declarations without type, function bodies
introduced by \verb|asm| with \verb|method_id| after modifiers;
\item function modifiers: \verb|impure|, \verb|inline|, \verb|inline_ref|,
\verb|method_id| in \verb|0.4.6| positions;
\item types: \verb|int| as an alias for \verb|i257|, and lowercase
\verb|cell|, \verb|slice|, \verb|builder|, \verb|cont|, \verb|tuple| as
aliases for the \verb|0.5.0| domain types;
\item statements: \verb|repeat Expr Block|, \verb|do Block until Expr|;
\item expressions: \verb|var ...| bindings, explicit type application
\verb|(int, slice) (n, s) = ...| on the left of \verb|=|, bare unit
return \verb|return (())|; method calls of the form
\verb|receiver~name(args)|.
\end{itemize}

Cross-edition interoperation:

\begin{itemize}
\item a \verb|0.5.0| file may \verb|use| items from a \verb|0.4-compat| file,
provided the imported items have types expressible in the \verb|0.5.0|
surface;
\item a \verb|0.4-compat| file may call \verb|0.5.0| items that do not use
features absent from the compat surface (notably: \verb|Result|, typed
\verb|entry|, traits, or enums);
\item a \verb|0.4-compat| file may \emph{not} define a \verb|0.5.0|
\verb|entry| function; entry points cross the compatibility boundary
exactly once, at the outer edition of the contract's root module;
\item the contract's root module's edition determines the emitted ABI.
\end{itemize}

\paragraph{Removal schedule.} \verb|0.4-compat| is a migration surface, not a
long-term supported edition. It is scheduled for removal in a future edition
\verb|0.7.0|. Contracts that have not migrated by then must freeze on the
last compiler release that still accepts it.

\section{Canonical Corpus}
\label{sec:canonical-corpus}
The 10 Canonical Decisions above are normative on paper; the corpus is how
they are validated against real contracts and how cross-compiler conformance
is checked. The corpus is defined in a sibling document,
\verb|doc/func_v0.5.0-corpus.md|, which ships as version-pinned artefacts
under \verb|crypto/smartcont/v050/|.

Every compliance-seeking compiler must produce byte-identical TVM code-cell
output for every corpus entry at its pinned version. Corpus v1 includes:

\begin{itemize}
\item C1 --- Simple wallet (\verb|wallet3|);
\item C2 --- Plugin wallet (\verb|wallet-v4|);
\item C3 --- Highload wallet v2;
\item C4 --- Jetton wallet (synthetic, from TEP-74);
\item C5 --- Multisig (sketch; full rewrite in corpus v2);
\item --- Elector (deferred to corpus v2).
\end{itemize}

The corpus is intentionally small in v1: enough to pin the language, not so
large that freezing it blocks edition maturation. Corpus versions are
monotonic --- v2 may add entries but may not redefine entries already fixed
in v1, short of a spec-revision-grade compatibility break.

\paragraph{Issues surfaced by corpus v1, now resolved in main spec v3.}
Every issue that corpus v1 flagged as open has been folded into the main
specification:

\begin{itemize}
\item helper-function state mutation --- resolved in
\S\ref{sec:mut-args} by the existing \verb|mut| parameter qualifier;
there is no \verb|\&mut self| syntax;
\item \verb|Vec<T>| at getter and ABI boundaries --- resolved in
\S\ref{sec:vec-type};
\item \verb|EitherRefOrInline<T>| per-field encoding --- resolved in
\S\ref{sec:either-ref-inline};
\item bounce-body width enforcement --- resolved in
\S\ref{sec:entry-decode} by the \texttt{BounceBodyTooWide} compile-time
check;
\item canonical \verb|std::| module layout --- resolved in
Appendix~\ref{app:stdlib-map} below.
\end{itemize}

\appendix
\section{Standard Library Module Map}
\label{app:stdlib-map}
The canonical \verb|std::| module tree below is part of the
edition-\verb|0.5.0| ABI. Every conforming compiler must accept these module
paths and export the listed names. Item signatures are authoritative in the
modules themselves (shipped alongside the reference compiler); this appendix
fixes the \emph{tree} so that \verb|use| paths across source bases are
compatible.

\begin{longtable}{|l|p{0.58\textwidth}|}
\hline
\textbf{Module path} & \textbf{Key exports} \\ \hline
\endfirsthead
\hline
\textbf{Module path} & \textbf{Key exports} \\ \hline
\endhead
\verb|std::prelude| & auto-imported into every source file:
\verb|Result|, \verb|Option|, \verb|Ok|, \verb|Err|, \verb|Some|,
\verb|None|; the \verb|?| lowering relies on these exact names \\ \hline
\verb|std::cell| & \verb|Cell<T>|, \verb|RawCell|, \verb|Builder|,
\verb|Slice|, \verb|RawSlice|, \verb|CellError|, \verb|ParseError|;
conversions \verb|to_cell|, \verb|from_cell|, \verb|load_any|,
\verb|store_any|; \verb|Cell::from_raw|, \verb|Cell::as_raw| \\ \hline
\verb|std::address| & \verb|Address|, \verb|StdAddress|,
\verb|parse_std_addr|, \verb|Address::derive|, address codecs \\ \hline
\verb|std::coins| & \verb|Coins|, \verb|Coins::zero|, \verb|Coins::is_zero|,
arithmetic, safe conversions \\ \hline
\verb|std::bytes| & \verb|Bytes|, \verb|LenPrefixedBytes|, helpers for snake
encoding and byte-alignment assertions \\ \hline
\verb|std::vec| & \verb|Vec<T>|, \verb|EitherRefOrInline<T>|; iteration
traits on \verb|Vec| \\ \hline
\verb|std::dict| & \verb|Dict<K, V>|, \verb|DictEntry<K, V>|, iteration via
\verb|iter_sorted|, \verb|Dict::pop_min|, \verb|Dict::contains|,
\verb|Dict::insert|, \verb|Dict::remove|, \verb|Dict::get| \\ \hline
\verb|std::crypto| & \verb|check_signature|, \verb|slice_hash|,
\verb|sha256|, \verb|keccak256|, curve primitives that expose TVM hashing
opcodes \\ \hline
\verb|std::chain| & \verb|now|, \verb|global_id|, \verb|accept_message|,
\verb|raw_reserve|, \verb|get_data|, \verb|set_data|, contract-balance
and block-context readers \\ \hline
\verb|std::message| & \verb|send|, \verb|send_raw|, \verb|SendOptions|,
\verb|SendMode|, \verb|BounceMode|, \verb|BodyLayout|,
\verb|InternalContext<T>|, \verb|ExternalContext|,
\verb|BounceContext<T>|, \verb|SendError| \\ \hline
\verb|std::iter| & \verb|Iter|, \verb|IntoIter| traits (see
\S\ref{sec:cd-expr}); stdlib iterator impls for \verb|Vec|, \verb|Dict|,
ranges, fixed arrays, uniform tuples \\ \hline
\verb|std::result| & re-exports from \verb|std::prelude|; additional
combinators \verb|map|, \verb|map_err|, \verb|and_then|,
\verb|or_else|, \verb|ok_or|, \verb|expect| \\ \hline
\verb|std::option| & re-exports from \verb|std::prelude|; additional
combinators \verb|map|, \verb|unwrap|, \verb|unwrap_or|, \verb|ok_or|,
\verb|is_some|, \verb|is_none| \\ \hline
\verb|std::intrinsics| & low-level \verb|extern "tvm"| declarations for
every TVM opcode that the stdlib surface wraps, available to
\verb|unsafe| code only \\ \hline
\end{longtable}

Module-path collision rules:

\begin{itemize}
\item third-party libraries may not publish into the \verb|std::*|
namespace;
\item contract code may publish its own \verb|mod foo::bar| tree freely;
\item the \verb|use| resolver searches the current crate first, then
\verb|std::*| as a reserved root, then external dependencies listed in
the build manifest;
\item \verb|use std::prelude::*;| is applied implicitly at the top of
every source file under edition \verb|0.5.0|.
\end{itemize}

A full item-level stdlib reference (every exported signature with
documentation) is scheduled as a sibling document,
\verb|doc/func_v0.5.0-stdlib-ref.md|, to be written alongside the reference
compiler implementation. The tree above is the frozen surface; the
signatures inside each module may still evolve up to edition freeze.

\end{document}
